% ================================================================
%  CAPITOLO 2 - INTRODUZIONE AD ARDUINO
% ================================================================
\chapter{Introduzione ad Arduino}

\section{La scheda Arduino Uno}
\subsection*{Panoramica dei componenti principali}
\subsection*{Il microcontrollore ATmega328P}
\subsection*{Versioni e differenze principali}

% ----------------------------------------------------------------
\section{Alimentazione della scheda}
% ----------------------------------------------------------------

Per funzionare correttamente, la scheda Arduino Uno necessita di un'alimentazione a corrente continua con una tensione compresa tra \SI{5}{\volt} e \SI{12}{\volt} (a seconda della modalità scelta). La scheda è progettata per essere alimentata in tre modi diversi, ciascuno con caratteristiche e limiti specifici.

\subsection*{Modalità di alimentazione: USB, Vin e alimentazione esterna}

Arduino Uno può ricevere energia elettrica attraverso tre differenti ingressi:

\begin{enumerate}
    \item \textbf{Porta USB (tipo B)}: quando la scheda è collegata al computer tramite cavo USB, riceve una tensione stabile di \SI{5}{\volt} direttamente dalla porta USB del computer. Questa è la modalità più semplice e sicura per alimentare la scheda durante la fase di sviluppo e caricamento degli sketch. La corrente disponibile è limitata a \SI{500}{\milli\ampere} dalle specifiche USB 2.0.

    \item \textbf{Connettore di alimentazione esterna (jack-barrel DC)}: è il connettore cilindrico di colore nero presente sulla scheda (diametro esterno \SI{5.5}{\milli\meter}, diametro interno \SI{2.1}{\milli\meter}, positivo al centro). Accetta tensioni continue comprese tra \SI{7}{\volt} e \SI{12}{\volt}. La tensione in ingresso viene poi ridotta a \SI{5}{\volt} dal regolatore di tensione a bordo.

    \item \textbf{Pin Vin}: è un pin di ingresso presente sul connettore \texttt{POWER} della scheda. Può essere utilizzato per fornire una tensione non regolata (sempre tra \SI{7}{\volt} e \SI{12}{\volt}) direttamente da una batteria o da un alimentatore esterno, bypassando il jack DC ma passando comunque attraverso il regolatore di tensione.
\end{enumerate}

\begin{center}
\begin{tabular}{>{\bfseries}l l l l}
    \toprule
    Modalità           & Tensione in ingresso  & Note \\
    \midrule
    Porta USB          & \SI{5}{\volt} (stabile) & Il computer limita la corrente \\
    Jack DC (consigliato) & \SIrange{7}{12}{\volt} & Tensione > \SI{12}{\volt} surriscalda il regolatore \\
    Pin Vin            & \SIrange{7}{12}{\volt} & Non protegge da inversione di polarità \\
    \bottomrule
\end{tabular}
\end{center}

\begin{tcolorbox}[colback=red!4, colframe=red!60, title=\textbf{Attenzione – Pin 5V come ingresso}]
    Il pin etichettato \texttt{5V} sul connettore \texttt{POWER} è un'\textbf{uscita} di tensione regolata a \SI{5}{\volt}, non un ingresso. Fornire tensione direttamente su questo pin è \textbf{sconsigliato} perché si bypassano le protezioni del regolatore e si rischia di danneggiare irreversibilmente il microcontrollore se la tensione non è perfettamente stabile a \SI{5}{\volt}.
\end{tcolorbox}

\subsection*{Il regolatore di tensione a bordo}

Quando la scheda viene alimentata tramite jack DC o pin Vin con una tensione superiore a \SI{5}{\volt}, entra in funzione il \textbf{regolatore di tensione lineare} integrato sulla scheda. Su Arduino Uno si tratta tipicamente di un componente siglato \textbf{NCP1117ST50T3G} o \textbf{AMS1117-5.0}, in grado di fornire una tensione stabile di \SI{5}{\volt} a partire da un ingresso compreso tra \SI{7}{\volt} e \SI{12}{\volt}.

Il principio di funzionamento di un regolatore lineare è semplice: la tensione in eccesso ($V_{in} - \SI{5}{\volt}$) viene \textit{dissipata} sotto forma di calore. Maggiore è la differenza tra tensione in ingresso e tensione di uscita, maggiore sarà il calore prodotto.

\begin{equation}
    P_{\text{dissipata}} = (V_{in} - \SI{5}{\volt}) \times I_{\text{assorbita}}
\end{equation}

Ad esempio, se alimentiamo la scheda a \SI{12}{\volt} e il circuito assorbe complessivamente \SI{200}{\milli\ampere}, il regolatore deve dissipare:

\[
    P = (\SI{12}{\volt} - \SI{5}{\volt}) \times \SI{0.2}{\ampere} = \SI{1.4}{\watt}
\]
Questo livello di potenza dissipata è già significativo per un piccolo componente senza dissipatore, e può portare a temperature elevate che riducono la vita utile del regolatore e della scheda stessa. Per questo motivo è fortemente consigliato non superare i \SI{12}{\volt} di tensione in ingresso quando si alimenta Arduino tramite jack DC o pin Vin.

\begin{tcolorbox}[colback=orange!4, colframe=orange!70, title=\textbf{Perché non superare i \SI{12}{\volt}?}]
    Il regolatore di tensione di Arduino Uno può teoricamente accettare tensioni fino a \SI{20}{\volt} (limite massimo assoluto del datasheet), ma nella pratica tensioni superiori a \SI{12}{\volt} causano un surriscaldamento eccessivo del componente, con conseguente riduzione della vita utile della scheda o spegnimento per protezione termica. Per alimentazioni superiori a \SI{12}{\volt} è preferibile utilizzare un regolatore switching esterno.
\end{tcolorbox}

\subsection*{Selezione automatica della sorgente di alimentazione}

Arduino Uno è dotato di un circuito di \textbf{selezione automatica della sorgente di alimentazione}, realizzato tramite un amplificatore operazionale (comparatore) e un MOSFET a canale P. Questo circuito decide quale delle due possibili fonti di energia utilizzare:

\begin{itemize}
    \item Se è presente tensione sul jack DC (o sul pin Vin) superiore a circa \SI{6.6}{\volt}, il circuito scollega automaticamente l'alimentazione proveniente dalla porta USB e utilizza la tensione esterna.
    \item Se la tensione esterna scende al di sotto di tale soglia (o è assente), il circuito commuta sull'alimentazione USB.
\end{itemize}

Questa commutazione avviene in modo completamente trasparente per l'utente e garantisce che la scheda possa essere alimentata in modo sicuro anche quando sono collegate entrambe le sorgenti (ad esempio durante il caricamento di uno sketch da computer mentre la scheda è già alimentata da una batteria esterna).

% ----------------------------------------------------------------
\section{I pin di alimentazione}
% ----------------------------------------------------------------

Sul connettore \texttt{POWER} di Arduino Uno sono disponibili diversi pin dedicati all'alimentazione del circuito. Questi pin possono essere utilizzati per prelevare tensione stabilizzata da destinare ai componenti esterni (sensori, LED, moduli) oppure per fornire alimentazione alla scheda stessa.

\subsection*{Pin 5V e 3.3V: caratteristiche e limiti di corrente}

I pin contrassegnati con \texttt{5V} e \texttt{3.3V} forniscono tensioni stabilizzate e possono essere utilizzati per alimentare circuiti esterni compatibili con tali livelli di tensione.

\begin{itemize}
    \item \textbf{Pin 5V}: fornisce una tensione stabilizzata di \SI{5}{\volt}. Questa tensione è generata dal regolatore lineare integrato sulla scheda (quando l'alimentazione proviene da jack DC o Vin) oppure proviene direttamente dalla porta USB (quando la scheda è alimentata via USB).

    \item \textbf{Pin 3.3V}: fornisce una tensione stabilizzata di \SI{3.3}{\volt}, generata da un secondo regolatore lineare (tipicamente un LP2985) che preleva l'ingresso dalla linea a \SI{5}{\volt}. 
\end{itemize}

\begin{center}
\begin{tabular}{>{\bfseries}l l l}
    \toprule
    Pin & Tensione & Corrente massima consigliata \\
    \midrule
    \texttt{5V}  & \SI{5}{\volt}  & \SI{500}{\milli\ampere} (USB) / \SI{800}{\milli\ampere} (esterna) \\
    \texttt{3.3V} & \SI{3.3}{\volt} & \SI{50}{\milli\ampere} \\
    \bottomrule
\end{tabular}
\end{center}

\begin{tcolorbox}[colback=red!4, colframe=red!60, title=\textbf{Attenzione – Non fornire tensione sui pin 5V o 3.3V}]
    I pin \texttt{5V} e \texttt{3.3V} sono \textbf{uscite} di tensione regolata. Non devono essere utilizzati come ingressi per alimentare la scheda. Fornire una tensione esterna su questi pin, specialmente se superiore al valore nominale, può danneggiare irreversibilmente il regolatore e il microcontrollore.
\end{tcolorbox}

\subsection*{Pin Vin: ingresso per alimentazione non regolata}

Il pin \texttt{Vin} (abbreviazione di \textit{Voltage input}) è un ingresso di alimentazione non regolata. Può essere utilizzato per fornire tensione alla scheda quando non si impiega il jack DC o la porta USB. Le caratteristiche principali sono:

\begin{itemize}
    \item \textbf{Tensione di ingresso consigliata}: \SIrange{7}{12}{\volt}. Tensioni inferiori a \SI{7}{\volt} potrebbero non garantire il corretto funzionamento del regolatore a \SI{5}{\volt}, causando instabilità o reset del microcontrollore.
    \item \textbf{Tensione massima assoluta}: \SI{20}{\volt}. Tuttavia, come discusso in precedenza, tensioni superiori a \SI{12}{\volt} provocano un'eccessiva dissipazione di potenza nel regolatore, con conseguente surriscaldamento.
    \item \textbf{Protezione da inversione di polarità}: il pin \texttt{Vin} \textbf{non} è protetto da un diodo contro l'inversione di polarità (a differenza del jack DC). Collegare accidentalmente il positivo della batteria a \texttt{GND} e il negativo a \texttt{Vin} può distruggere istantaneamente il regolatore e altri componenti.
\end{itemize}

\subsection*{Pin GND: massa comune del circuito}

I pin etichettati \texttt{GND} (dall'inglese \textit{ground}, massa) costituiscono il riferimento comune a \SI{0}{\volt} per l'intero circuito. Su Arduino Uno sono presenti tre pin \texttt{GND} sul connettore \texttt{POWER} e altri due pin \texttt{GND} sui connettori digitali (vicino al pin 13 e al pin \texttt{AREF}). Tutti i pin \texttt{GND} sono collegati internamente alla stessa massa.

\begin{tcolorbox}[colback=blue!4, colframe=blue!60, title=\textbf{Buona pratica – Masse comuni}]
    Quando si utilizza un'alimentazione esterna per pilotare carichi di potenza (motori, relè, strisce LED), è \textbf{obbligatorio} collegare la massa dell'alimentazione esterna a un pin \texttt{GND} di Arduino. Questo garantisce che i segnali di controllo inviati da Arduino abbiano lo stesso riferimento di tensione del circuito di potenza, evitando malfunzionamenti o letture errate.
\end{tcolorbox}

\subsection*{Corrente massima erogabile e protezioni}

La scheda Arduino Uno incorpora alcune protezioni per limitare i danni in caso di utilizzo improprio, ma non è progettata per resistere a condizioni di sovraccarico prolungato.

\begin{itemize}
    \item \textbf{Fusibile ripristinabile}: sulla linea di alimentazione a \SI{5}{\volt} proveniente dalla porta USB è presente un fusibile a coefficiente di temperatura positivo (PTC). Questo componente aumenta la propria resistenza quando la corrente supera una soglia (tipicamente \SI{500}{\milli\ampere}), limitando la corrente e proteggendo la porta USB del computer. Quando il sovraccarico cessa, il componente si raffredda e ripristina la conduzione normale. Il fusibile \textbf{non} protegge il regolatore a \SI{5}{\volt} quando la scheda è alimentata tramite jack DC o Vin.

    \item \textbf{Diodo di protezione sul jack DC}: un diodo in serie al jack DC impedisce che una tensione inversa applicata accidentalmente raggiunga il regolatore. Tuttavia, questo diodo introduce una caduta di tensione di circa \SI{0.7}{\volt}, che deve essere considerata nella scelta dell'alimentatore.

    \item \textbf{Diodi di clamping sui pin di I/O}: ciascun pin digitale è protetto da diodi interni che limitano la tensione a un massimo di $V_{cc} + \SI{0.5}{\volt}$ e a un minimo di $\SI{-0.5}{\volt}$. Questi diodi proteggono il microcontrollore da scariche elettrostatiche (ESD) e da brevi sovratensioni, ma non sono in grado di dissipare potenze significative. Applicare tensioni negative o superiori a \SI{5.5}{\volt} per periodi prolungati può distruggere i diodi e il pin corrispondente.
\end{itemize}

\begin{tcolorbox}[colback=red!4, colframe=red!60, title=\textbf{Riepilogo – Limiti operativi assoluti}]
    \begin{center}
    \begin{tabular}{>{\bfseries}l l l}
        Parametro & $\quad$ & \textbf{Limite} \vspace{0.1cm}\\
        $V_\texttt{MAX}$ su jack DC / Vin & & \SI{20}{\volt} (assoluto), \textbf{\SI{12}{\volt} raccomandato} \\
        $V_\texttt{MIN}$ su jack DC / Vin  & & \SI{7}{\volt} \\
        $I_\texttt{MAX}$ da pin \texttt{5V} && \SI{800}{\milli\ampere} (dipende da Vin e temperatura) \\
        $I_\texttt{MAX}$ da pin \texttt{3.3V} && \SI{50}{\milli\ampere} \\
        $I_\texttt{MAX}$ totale da tutti i pin di I/O && \SI{200}{\milli\ampere} \\
    \end{tabular}
    \end{center}
\end{tcolorbox}

% ----------------------------------------------------------------
\section{I pin digitali}
% ----------------------------------------------------------------

I pin digitali di Arduino Uno sono i connettori contrassegnati dai numeri da 0 a 13 sul bordo superiore e inferiore della scheda (Figura ref pinout). Ciascuno di questi pin può essere configurato via software per comportarsi come \textbf{ingresso} (INPUT) o come \textbf{uscita} (OUTPUT). Quando configurato come uscita, il pin può fornire una tensione di \SI{5}{\volt} (livello logico HIGH) o \SI{0}{\volt} (livello logico LOW). Quando configurato come ingresso, il pin è in grado di leggere la tensione presente sul terminale e interpretarla come un segnale logico HIGH o LOW.

\subsection*{Configurazione come ingresso o uscita}

La direzione del flusso del segnale è definita via software, ma corrisponde a una precisa configurazione hardware interna al microcontrollore. Ogni pin digitale è collegato a due sottocircuiti principali: uno di uscita (driver) e uno di ingresso (buffer di lettura).

\begin{enumerate}
    \item \textbf{Configurazione come uscita (\texttt{OUTPUT})}: viene attivato il circuito di uscita, costituito da una coppia di transistor che collegano il pin alla tensione di alimentazione oppure a massa. Il microcontrollore impone quindi direttamente il livello logico del pin (HIGH $\approx$ \SI{5}{\volt}, LOW $\approx$ \SI{0}{\volt}) ed è in grado di fornire o assorbire corrente verso il circuito esterno.

    \item \textbf{Configurazione come ingresso (\texttt{INPUT})}: il circuito di uscita viene disattivato, quindi il microcontrollore non impone alcun livello di tensione sul pin. Viene invece attivato il circuito di ingresso, che misura la tensione presente sul pin senza influenzarla. L'impedenza vista dall'esterno risulta molto elevata e la corrente assorbita è trascurabile.
\end{enumerate}

Quando un pin è configurato come ingresso, il suo livello logico dipende esclusivamente dal circuito esterno. Se il pin non è collegato a una sorgente definita, la tensione può assumere valori imprevedibili a causa dei disturbi elettrici.



\begin{comment}
\begin{figure}[h]
    \centering
    \includegraphics[width=0.8\textwidth]{img/Schema_Pin_GPIO.png}
    \caption{Struttura semplificata di un pin di I/O del microcontrollore ATmega328P.}
    \label{fig:struttura_pin_gpio}
\end{figure}
\end{comment}

\subsection*{Livelli di tensione e soglie logiche}

Quando un pin digitale è configurato come \textbf{uscita}, i livelli di tensione sono ben definiti: \SI{0}{\volt} per LOW e \SI{5}{\volt} per HIGH (in realtà circa \SI{4.8}{\volt}–\SI{4.9}{\volt} a causa della caduta interna sul transistor di uscita).

Quando il pin è configurato come \textbf{ingresso}, la tensione applicata dall'esterno viene interpretata secondo le \textbf{soglie logiche} definite dal costruttore del microcontrollore ATmega328P:

\begin{center}
\begin{tabular}{>{\bfseries}l l l}
    \toprule
    Livello logico & Tensione in ingresso ($V_{in}$) & Interpretazione \\
    \midrule
    LOW  & $V_{in} < \SI{1.5}{\volt}$ & Il pin legge 0 logico \\
    HIGH & $V_{in} > \SI{3.0}{\volt}$ & Il pin legge 1 logico \\
    \bottomrule
\end{tabular}
\end{center}

\begin{tcolorbox}[colback=red!4, colframe=red!60, title=\textbf{Attenzione – La zona proibita}]
    Qualsiasi tensione compresa tra \SI{1.5}{\volt} e \SI{3.0}{\volt} cade in una \textbf{zona indeterminata}: il microcontrollore potrebbe interpretare il segnale in modo casuale come HIGH o LOW, generando comportamenti imprevedibili. Per questo motivo è necessario assicurarsi che i segnali digitali in ingresso abbiano transizioni nette tra i due stati logici.
\end{tcolorbox}

\subsection*{Corrente massima per pin e per l'intera scheda}

I pin di Arduino non sono in grado di erogare o assorbire una quantità illimitata di corrente. Dal datasheet dell'ATmega328P e dalle specifiche della scheda Arduino Uno si ricavano i seguenti limiti operativi:

\begin{center}
\begin{tabular}{>{\bfseries}l l}
    \toprule
    Parametro & Valore massimo \\
    \midrule
    Corrente erogata da un singolo pin configurato come uscita HIGH & \SI{40}{\milli\ampere} \\
    Corrente assorbita da un singolo pin configurato come uscita LOW  & \SI{40}{\milli\ampere} \\
    Corrente totale erogata dall'insieme di tutti i pin di uscita     & \SI{200}{\milli\ampere} \\
    \bottomrule
\end{tabular}
\end{center}

\begin{tcolorbox}[colback=red!4, colframe=red!60, title=\textbf{Attenzione – Superamento dei limiti di corrente}]
    Superare anche per breve tempo i \SI{40}{\milli\ampere} su un singolo pin o i \SI{200}{\milli\ampere} complessivi può danneggiare irreversibilmente il microcontrollore o il regolatore di tensione della scheda. Per pilotare carichi che richiedono correnti superiori (motori, relè, strisce LED ad alta luminosità) è \textbf{obbligatorio} utilizzare componenti esterni come transistor o circuiti integrati driver.
\end{tcolorbox}

\subsection*{Resistenze di pull-up interne}

Quando un pin digitale è configurato come \textbf{ingresso} e non è collegato ad alcun circuito esterno (né a \SI{5}{\volt} né a massa), la sua tensione non è definita e può variare in modo imprevedibile a causa dei campi elettromagnetici ambientali e dei disturbi elettrici. Questa condizione è nota come \textbf{pin flottante} e porta a letture casuali e inaffidabili.

Per evitare questo problema, il microcontrollore dispone di \textbf{resistenze di pull-up interne} attivabili via software. Una resistenza di pull-up è un resistore di elevato valore (tipicamente \SIrange{20}{50}{\kilo\ohm}) che collega internamente il pin alla tensione di alimentazione \SI{5}{\volt}.

Quando la resistenza di pull-up è attiva, in assenza di un circuito esterno che imponga un diverso livello di tensione, il pin assume un livello logico HIGH. Se invece il pin viene collegato a massa attraverso un circuito esterno a bassa impedenza, la tensione scende a \SI{0}{\volt} (livello LOW), poiché la corrente che attraversa la resistenza di pull-up è molto piccola e non è sufficiente a mantenere il livello alto.

In questo modo si garantisce che il pin presenti sempre un livello logico definito, evitando condizioni di indeterminazione.


\begin{tcolorbox}[colback=blue!4, colframe=blue!60, title=\textbf{Buona pratica – Evitare pin flottanti}]
    Ogni pin di ingresso \textbf{deve sempre} essere portato a un livello di tensione definito, o tramite un collegamento esterno (ad esempio con una resistenza di pull-down o pull-up) o attivando le resistenze di pull-up interne del microcontrollore. Non lasciare mai pin di ingresso scollegati.
\end{tcolorbox}

% ----------------------------------------------------------------
\section{I pin analogici}
% ----------------------------------------------------------------

I pin analogici di Arduino Uno sono sei, contrassegnati dalle etichette \texttt{A0} – \texttt{A5} sul lato sinistro della scheda. A differenza dei pin digitali, che possono assumere solo due stati logici (\texttt{HIGH} o \texttt{LOW}), i pin analogici sono in grado di misurare una \textbf{gamma continua di tensioni} compresa tra \SI{0}{\volt} e un valore massimo di riferimento (tipicamente \SI{5}{\volt}). Questa capacità è resa possibile dal \textbf{convertitore analogico-digitale} (ADC) integrato nel microcontrollore ATmega328P.

\subsection*{Il convertitore analogico-digitale (ADC) a 10 bit}

Un convertitore analogico-digitale è un circuito elettronico che trasforma una tensione continua (segnale analogico) in un numero digitale (segnale discreto) che può essere elaborato dal microcontrollore. L'ADC dell'ATmega328P ha una risoluzione di \textbf{10 bit}. Ciò significa che il valore letto viene rappresentato con un numero intero a 10 bit, che può assumere $2^{10} = 1024$ valori distinti, da \textbf{0} a \textbf{1023}.

La relazione tra la tensione in ingresso $V_{in}$ e il valore numerico restituito dall'ADC è lineare:

\begin{equation}
    \text{valore ADC} = \frac{V_{in}}{V_{ref}} \times 1023
\end{equation}

dove $V_{ref}$ è la \textbf{tensione di riferimento} dell'ADC. Conoscendo il valore letto e la tensione di riferimento, è possibile risalire alla tensione originale:

\begin{equation}
    V_{in} = \frac{\text{valore ADC}}{1023} \times V_{ref}
\end{equation}

\begin{tcolorbox}[colback=blue!4, colframe=blue!60, title=\textbf{Nota – Il valore 1023}]
    Poiché i valori vanno da 0 a 1023, il massimo valore numerico corrisponde a $V_{in} = V_{ref}$. Alcune formule approssimano dividendo per 1024 invece di 1023, ma la differenza è trascurabile nella maggior parte delle applicazioni pratiche. La documentazione ufficiale di Arduino utilizza il fattore $1023$.
\end{tcolorbox}

\subsection*{Tensione di riferimento e risoluzione}

La \textbf{tensione di riferimento} ($V_{ref}$) determina il fondo scala dell'ADC, ovvero la massima tensione che può essere misurata. Per impostazione predefinita, Arduino Uno utilizza la propria tensione di alimentazione come riferimento, ovvero \SI{5}{\volt}. In questa configurazione, l'ADC può misurare tensioni comprese tra \SI{0}{\volt} e \SI{5}{\volt}.

Con $V_{ref} = \SI{5}{\volt}$, la \textbf{risoluzione} dell'ADC — ovvero la minima variazione di tensione che il convertitore è in grado di distinguere — è data da:

\begin{equation}
    \text{Risoluzione} = \frac{V_{ref}}{1024} \approx \frac{\SI{5}{\volt}}{1024} \approx \SI{4.9}{\milli\volt}
\end{equation}

Ciò significa che variazioni di tensione inferiori a circa \SI{5}{\milli\volt} non vengono rilevate, e il valore letto rimane invariato.

È possibile modificare la tensione di riferimento utilizzando il pin \texttt{AREF} (Analog Reference) o selezionando riferimenti interni alternativi tramite software.

\begin{tcolorbox}[colback=red!4, colframe=red!60, title=\textbf{Attenzione – Utilizzo di AREF}]
    Se si applica una tensione di riferimento esterna sul pin \texttt{AREF}, è \textbf{obbligatorio} configurare il software per utilizzare il riferimento esterno \textbf{prima} di effettuare qualsiasi lettura analogica. In caso contrario, si rischia di danneggiare il microcontrollore a causa di un conflitto tra il riferimento interno e quello esterno.
\end{tcolorbox}



\begin{tcolorbox}[colback=orange!4, colframe=orange!70, title=\textbf{Tempo di acquisizione e frequenza di campionamento}]
    L'ADC dell'ATmega328P richiede un certo tempo per stabilizzare la lettura dopo una variazione del segnale di ingresso. La frequenza massima di campionamento è di circa \SI{15}{\kilo\hertz} (circa \num{15000} letture al secondo) nella configurazione predefinita. Letture più rapide possono essere ottenute modificando i registri interni, ma a scapito della precisione.
\end{tcolorbox}

\subsection*{Utilizzo come pin digitali aggiuntivi}

I pin analogici \texttt{A0} – \texttt{A5} possono essere utilizzati anche come \textbf{pin digitali di I/O} aggiuntivi. In questa modalità, si comportano esattamente come i pin digitali standard (0–13), supportando le funzioni \texttt{pinMode()}, \texttt{digitalWrite()} e \texttt{digitalRead()}. Per riferirsi a un pin analogico in modalità digitale, si può utilizzare l'etichetta corrispondente (es. \texttt{A0}) oppure il numero equivalente (per Arduino Uno, \texttt{A0} corrisponde al pin 14, \texttt{A1} al 15, e così via fino a \texttt{A5} che corrisponde al pin 19).

Questa caratteristica è particolarmente utile quando i pin digitali standard sono già tutti occupati e si necessita di qualche I/O digitale supplementare.

\begin{tcolorbox}[colback=blue!4, colframe=blue!60, title=\textbf{Nota – PWM sui pin analogici}]
    A differenza di alcuni pin digitali (3, 5, 6, 9, 10, 11), i pin analogici \textbf{non} supportano l'uscita PWM (\texttt{analogWrite()}). Tentare di utilizzare \texttt{analogWrite()} su un pin analogico non produrrà l'effetto desiderato.
\end{tcolorbox}

\begin{center}
\begin{tabular}{>{\bfseries}l l l}
    \toprule
    Pin analogico & Equivalente digitale & Supporta PWM? \\
    \midrule
    A0 & 14 & No \\
    A1 & 15 & No \\
    A2 & 16 & No \\
    A3 & 17 & No \\
    A4 & 18 & No (utilizzato anche per I²C SDA) \\
    A5 & 19 & No (utilizzato anche per I²C SCL) \\
    \bottomrule
\end{tabular}
\end{center}

\begin{tcolorbox}[colback=red!4, colframe=red!60, title=\textbf{Attenzione – Pin A4 e A5 e comunicazione I²C}]
    I pin analogici \texttt{A4} (SDA) e \texttt{A5} (SCL) sono condivisi con l'interfaccia di comunicazione I²C (TWI). Se si utilizzano dispositivi I²C (come display LCD o sensori), questi pin non sono disponibili per l'uso come I/O digitale generico. In caso di conflitto, è possibile utilizzare pin analogici differenti o disabilitare l'I²C via software.
\end{tcolorbox}

% ----------------------------------------------------------------
\section{Pin di comunicazione seriale (TX e RX)}
% ----------------------------------------------------------------

La comunicazione seriale è uno dei metodi più semplici e diffusi per trasmettere dati tra due dispositivi. Arduino Uno dispone di una porta seriale hardware integrata nel microcontrollore, accessibile attraverso i pin digitali 0 e 1, contrassegnati rispettivamente con le etichette \texttt{RX} (Receive) e \texttt{TX} (Transmit).

\subsection*{La porta seriale hardware USART}

All'interno del microcontrollore ATmega328P è presente un modulo hardware dedicato denominato \textbf{USART} (\textit{Universal Synchronous and Asynchronous serial Receiver and Transmitter}). Questo circuito è in grado di gestire la trasmissione e la ricezione di dati seriali in modo autonomo, senza impegnare la CPU per ogni singolo bit trasmesso.

Nella modalità più comune, detta \textbf{asincrona}, la comunicazione avviene secondo il seguente schema:

\begin{itemize}
    \item I dati vengono trasmessi un bit alla volta su un singolo filo (più un filo di massa come riferimento).
    \item Il trasmettitore e il ricevitore devono concordare preventivamente una \textbf{velocità di trasmissione} (\textit{baud rate}), espressa in bit al secondo (bps).
    \item Ogni byte di dati è incorniciato da un \textbf{bit di start} (sempre a livello logico LOW) e da uno o più \textbf{bit di stop} (sempre a livello logico HIGH).
    \item Opzionalmente può essere aggiunto un \textbf{bit di parità} per il rilevamento di errori di trasmissione.
\end{itemize}

Su Arduino Uno, la USART è tipicamente c onfigurata con un baud rate di \SI{9600}{baud} (9600 bit al secondo) e utilizza il formato \textbf{8-N-1}: 8 bit di dati, nessun bit di parità, 1 bit di stop.

\begin{center}
\begin{tabular}{>{\bfseries}l l}
    \toprule
    Parametro & Valore tipico su Arduino Uno \\
    \midrule
    Baud rate predefinito & \SI{9600}{\texttt{baud}} (configurabile fino a \SI{115200}{\texttt{baud}}) \\
    Bit di dati & 8 \\
    Bit di parità & Nessuno (N) \\
    Bit di stop & 1 \\
    Livelli di tensione & \SI{0}{\volt} (LOW) e \SI{5}{\volt} (HIGH) \\
    \bottomrule
\end{tabular}
\end{center}

\begin{tcolorbox}[colback=blue!4, colframe=blue!60, title=\textbf{Approfondimento – Differenza tra USART e UART}]
    Il termine \textbf{UART} (\textit{Universal Asynchronous Receiver-Transmitter}) indica un modulo che gestisce esclusivamente la comunicazione asincrona. Il termine \textbf{USART} include anche la modalità sincrona, in cui viene trasmesso un segnale di clock separato. L'ATmega328P dispone di una USART, ma su Arduino viene utilizzata quasi esclusivamente in modalità asincrona (UART).
\end{tcolorbox}

\subsection*{Pin 0 (RX) e pin 1 (TX)}

I pin digitali \textbf{0 (RX)} e \textbf{1 (TX)} sono fisicamente collegati alla USART del microcontrollore e svolgono funzioni complementari:

\begin{itemize}
    \item \textbf{Pin 0 – RX (Receive)}: è l'ingresso della porta seriale. Riceve i dati provenienti da un altro dispositivo (ad esempio un computer, un modulo GPS, un altro Arduino) e li rende disponibili al microcontrollore per la lettura.
    \item \textbf{Pin 1 – TX (Transmit)}: è l'uscita della porta seriale. Trasmette i dati generati da Arduino verso un dispositivo esterno.
\end{itemize}

Quando si utilizza la comunicazione seriale tra due dispositivi, il pin \texttt{TX} di un dispositivo deve essere collegato al pin \texttt{RX} dell'altro, e viceversa. I due dispositivi devono inoltre condividere un riferimento di massa comune (\texttt{GND}).



È importante notare che i pin 0 e 1 sono \textbf{condivisi} tra la comunicazione seriale con il computer (attraverso il chip di conversione USB-Seriale) e qualsiasi circuito esterno collegato a questi pin. Questa condivisione ha implicazioni pratiche rilevanti, discusse nella sottosezione seguente.

\subsection*{Interferenza durante il caricamento degli sketch}

Quando si carica un nuovo programma su Arduino tramite la porta USB, i dati vengono inviati dal computer alla scheda attraverso la connessione seriale. Durante questa fase, il chip di conversione USB-Seriale (un ATmega16U2 su Arduino Uno Rev3) comunica direttamente con i pin \texttt{RX} e \texttt{TX} del microcontrollore ATmega328P.

Se in quel momento sono collegati componenti esterni ai pin 0 e 1, si possono verificare due tipi di interferenza:

\begin{enumerate}
    \item \textbf{Disturbo del segnale di caricamento}: un circuito esterno collegato al pin \texttt{RX} (pin 0) potrebbe imporre un livello di tensione che interferisce con i dati in arrivo dal computer, causando errori o l'impossibilità di caricare lo sketch.
    \item \textbf{Danneggiamento dei componenti esterni}: durante la programmazione, il pin \texttt{TX} (pin 1) trasmette dati verso il computer. Se a questo pin è collegato un componente sensibile (ad esempio un LED senza resistenza adeguata, o l'ingresso di un altro circuito a bassa tensione), la corrente di trasmissione potrebbe danneggiarlo.
\end{enumerate}

\begin{tcolorbox}[colback=red!4, colframe=red!60, title=\textbf{Attenzione – Scollegare i circuiti esterni prima del caricamento}]
    È buona norma \textbf{scollegare temporaneamente} qualsiasi circuito collegato ai pin 0 e 1 prima di caricare un nuovo sketch su Arduino. In alternativa, se il circuito esterno deve rimanere connesso, è possibile utilizzare pin seriali software (attraverso la libreria \texttt{SoftwareSerial}) su altri pin digitali, lasciando i pin 0 e 1 dedicati esclusivamente alla comunicazione con il computer.
\end{tcolorbox}

Una volta completato il caricamento, la comunicazione seriale tra Arduino e il computer può proseguire normalmente per lo scambio di dati a runtime (ad esempio per visualizzare i valori letti da un sensore sul \textit{Serial Monitor}).

\subsection*{La comunicazione con il computer tramite USB}

Sebbene i pin 0 e 1 siano fisicamente i terminali della comunicazione seriale, l'utente interagisce con essi principalmente attraverso la \textbf{porta USB} della scheda. Questa apparente contraddizione è risolta dal \textbf{chip di conversione USB-Seriale} presente sulla scheda.

Su Arduino Uno Rev3, questo ruolo è svolto da un secondo microcontrollore, l'\textbf{ATmega16U2}, programmato per comportarsi come un ponte tra il protocollo USB (utilizzato dal computer) e il protocollo seriale TTL (utilizzato dall'ATmega328P). Il flusso dei dati è il seguente:

\begin{enumerate}
    \item Il computer invia dati attraverso la porta USB emulando una porta seriale virtuale (COM su Windows, \texttt{/dev/tty*} su Linux e macOS).
    \item Il chip ATmega16U2 riceve i dati via USB, li converte in segnali seriali TTL e li inoltra al pin \texttt{RX} (pin 0) dell'ATmega328P.
    \item L'ATmega328P legge i dati dal pin \texttt{RX} e, se necessario, risponde trasmettendo dati sul pin \texttt{TX} (pin 1).
    \item L'ATmega16U2 riceve i dati dal pin \texttt{TX}, li converte in pacchetti USB e li invia al computer.
\end{enumerate}

Questo meccanismo è completamente trasparente per l'utente, che vede la scheda Arduino come una periferica seriale collegata al computer.

\begin{tcolorbox}[colback=blue!4, colframe=blue!60, title=\textbf{Nota – I LED TX e RX sulla scheda}]
    Sulla scheda Arduino Uno sono presenti due piccoli LED contrassegnati con \texttt{TX} e \texttt{RX}. Questi LED lampeggiano ogni volta che viene trasmesso o ricevuto un dato attraverso la porta seriale USB. Il loro lampeggiamento è particolarmente utile per verificare visivamente che la comunicazione sia in corso, ad esempio durante il caricamento di uno sketch o durante lo scambio di dati con il \textit{Serial Monitor}.
\end{tcolorbox}

La comunicazione seriale via USB rappresenta lo strumento principale per il \textit{debugging} dei programmi su Arduino: attraverso le funzioni \texttt{Serial.print()} e \texttt{Serial.println()} è possibile inviare messaggi e valori di variabili al computer, visualizzandoli nell'ambiente di sviluppo integrato.

\section{Pin di interrupt}
\subsection*{Concetto di interrupt hardware}
\subsection*{Pin 2 e 3: interrupt esterni disponibili}
\subsection*{Utilizzo tipico e vantaggi}

\section{L'ambiente di sviluppo Arduino IDE}
\subsection*{Installazione e configurazione}
\subsection*{Interfaccia principale: barra degli strumenti, editor e console}
\subsection*{Selezione della scheda e della porta seriale}
\subsection*{Verifica e caricamento di uno sketch}

\section{Struttura fondamentale di un programma}
\subsection*{La funzione \texttt{setup()}: inizializzazione}
\subsection*{La funzione \texttt{loop()}: esecuzione ciclica}
\subsection*{Il ruolo delle librerie}